\section{Related Work}
\label{sec:related}

\paragraph{Querying compressed data.}
Executing directly on compressed representations is a long-running theme in data
management. Column stores compress each column with a lightweight scheme and,
crucially, evaluate predicates on the encoded form rather than decoding
first~\cite{abadi2006integrating,boncz2005monetdb,zukowski2006super}; the same
argument drives modern lakehouse formats~\cite{melnik2010dremel,kuschewski2023btrblocks}
and fast integer codecs~\cite{lemire2015decoding}. Succinct~\cite{agarwal2015succinct}
pushes further, storing a compressed representation that supports search and random
access without expansion, building on compressed self-indexes~\cite{ferragina2000opportunistic,gagie2020fully}
and the broader compact-data-structures program~\cite{jacobson1989space,navarro2016compact}.
Our contribution is not the idea of computing on compressed data but a
characterization of \emph{which} compressed formats admit it
(\S\ref{sec:principles}), and a demonstration on a scientific interchange format
whose expanded form is the thing that no longer fits. Where column stores decode a
column to operate on it and self-indexes optimize for locate-style queries, we keep
a grammar-compressed representation resident and fold over it in place, which suits
dense scans that touch every element exactly once.

\paragraph{Grammar compression.}
Replacing repeated substrings with grammar rules is
classical~\cite{ziv1977universal,nevill1997identifying,larsson2000offline}, and the
observation that a straight-line grammar can be traversed without full expansion
underlies work on compressed pattern matching. The container we build
on~\cite{gfaz2026} applies an iterative 2-mer grammar to pangenome traversals;
sqz~\cite{heringer2025human} applies heap-based byte-pair encoding to the same
records. What this paper adds is that the grammar is a usable execution substrate
for whole analyses, not only for string queries, once it is paired with namespace
partitioning, record independence, and orientation closure.

\paragraph{Genomic storage formats.}
Genomics has repeatedly solved storage by designing a format rather than by
applying a general compressor: BAM and CRAM for aligned
reads~\cite{li2009sequence,bonfield2022cram}, block-gzip with an index for random
access into flat records~\cite{li2011tabix}, and specialized codecs for
sequence collections~\cite{yang2024gpufastqlz}. These are archival and
retrieval formats: they optimize for storing and fetching regions, and an analysis
still runs on decoded records. GBZ~\cite{siren2022gbz}, built on the
GBWT~\cite{siren2020haplotype,siren2021pangenomics}, is the exception and the
closest relative of this work, in that it supports navigation on the compressed
haplotype collection directly; \S\ref{subsec:bg-alternatives} explains why an index
tuned for locate-style navigation is a different structure from the one a dense
fold wants, and \S\ref{subsec:eval-gbz} measures the difference.

\paragraph{Pangenome analysis tools.}
The tools we compare against are the field's standard implementations, and each is
materialization-based by design: \cmd{vg} loads a graph and a snarl
decomposition~\cite{garrison2018variation,hickey2020genotyping,paten2018superbubbles},
\cmd{odgi} builds a succinct whole-graph object~\cite{guarracino2022odgi}, and
Panacus parses records into item tables~\cite{parmigiani2024panacus}. Graph
construction pipelines~\cite{hickey2024pangenome,garrison2024building} sit upstream
of all of this and are unaffected by our work. Our \cmd{deconstruct} algorithm is,
to our knowledge, the first to derive a reference-relative VCF without a
materialized graph or a general snarl index, by separating topology-only site
discovery from a streaming allele fold (\S\ref{sec:deconstruct}).
